\chapter{שיטות אימון ואופטימיזציה}

אימון מודלים של למידה עמוקה מצריך שיטות אופטימיזציה מתוחכמות. פרק זה סוקר את
האלגוריתמים המרכזיים לאימון רשתות נוירונים, כולל שיטות ירידת גרדיאנט, רגולריזציה,
ואסטרטגיות תזמון קצב למידה.

\section{ירידת גרדיאנט}

אלגוריתם ירידת הגרדיאנט (\textenglish{Gradient Descent}) הוא הבסיס לאימון
רשתות נוירונים. עדכון הפרמטרים מבוצע לפי:

\begin{equation}
    \theta_{t+1} = \theta_t - \eta \nabla_\theta \mathcal{L}(\theta_t)
\end{equation}

כאשר \(\eta\) הוא קצב הלמידה (\textenglish{Learning Rate}) ו-\(\mathcal{L}\)
היא פונקציית ההפסד.

\subsection{ירידת גרדיאנט סטוכסטית}

ירידת גרדיאנט סטוכסטית (\textenglish{Stochastic Gradient Descent},
קיצור: \textenglish{SGD}) מעדכנת את הפרמטרים על בסיס מדגם אקראי בלבד:

\begin{equation}
    \theta_{t+1} = \theta_t - \eta \nabla_\theta \mathcal{L}(\theta_t;\, x^{(i)},\, y^{(i)})
\end{equation}

\subsection{מומנטום}

שיטת המומנטום (\textenglish{Momentum}) מוסיפה מהירות להתכנסות:

\begin{equation}
    v_{t+1} = \beta v_t + (1 - \beta) \nabla_\theta \mathcal{L}(\theta_t)
\end{equation}

\begin{equation}
    \theta_{t+1} = \theta_t - \eta v_{t+1}
\end{equation}

\section{אופטימייזרים מתקדמים}

\subsection{\textenglish{Adam}}

אלגוריתם \textenglish{Adam} (\textenglish{Adaptive Moment Estimation})
הוא הנפוץ ביותר בלמידה עמוקה מודרנית. הוא משלב מומנטום ראשון ושני:

\begin{equation}
    m_t = \beta_1 m_{t-1} + (1 - \beta_1) g_t, \qquad
    v_t = \beta_2 v_{t-1} + (1 - \beta_2) g_t^2
\end{equation}

\begin{equation}
    \hat{m}_t = \frac{m_t}{1 - \beta_1^t}, \qquad
    \hat{v}_t = \frac{v_t}{1 - \beta_2^t}
\end{equation}

\begin{equation}
    \theta_{t+1} = \theta_t - \frac{\eta}{\sqrt{\hat{v}_t} + \epsilon} \hat{m}_t
\end{equation}

הערכים המומלצים הם \(\beta_1 = 0.9\), \(\beta_2 = 0.999\), ו-\(\epsilon = 10^{-8}\).

\subsection{\textenglish{AdamW}}

\textenglish{AdamW} מוסיף ריקבון משקלות (\textenglish{Weight Decay}) כרגולריזציה:

\begin{equation}
    \theta_{t+1} = \theta_t - \eta \left( \frac{\hat{m}_t}{\sqrt{\hat{v}_t} + \epsilon}
    + \lambda \theta_t \right)
\end{equation}

\section{רגולריזציה}

\subsection{\textenglish{Dropout}}

\textenglish{Dropout} הוא טכניקת רגולריזציה שמכבה נוירונים אקראית באימון:

\begin{equation}
    \tilde{h}_i = \frac{r_i}{1 - p} \cdot h_i, \qquad r_i \sim \text{Bernoulli}(1 - p)
\end{equation}

כאשר \(p\) הוא הסתברות הכיבוי (\textenglish{Drop Rate}).

\subsection{נרמול שכבה}

נרמול שכבה (\textenglish{Layer Normalization}) משמש נרחב בארכיטקטורת
\textenglish{Transformer}:

\begin{equation}
    \hat{x}_i = \frac{x_i - \mu}{\sqrt{\sigma^2 + \epsilon}}, \qquad
    y_i = \gamma \hat{x}_i + \beta
\end{equation}

\section{לוח זמנים לקצב למידה}

לוח זמנים חמים (\textenglish{Warmup Schedule}) הוא נפוץ באימון
\textenglish{Transformer}:

\begin{equation}
    \eta_t = d_{\text{model}}^{-0.5} \cdot
    \min\!\left( t^{-0.5},\; t \cdot t_{\text{warmup}}^{-1.5} \right)
\end{equation}

\section{דוגמת קוד: לולאת אימון}

\begin{LTR}
\begin{verbatim}
import torch
from torch.optim import AdamW
from torch.optim.lr_scheduler import CosineAnnealingLR

model = TransformerModel(...)
optimizer = AdamW(model.parameters(), lr=3e-4, weight_decay=0.01)
scheduler = CosineAnnealingLR(optimizer, T_max=num_epochs)

for epoch in range(num_epochs):
    for batch in dataloader:
        optimizer.zero_grad()
        inputs, targets = batch
        logits = model(inputs)
        loss = criterion(logits, targets)
        loss.backward()
        torch.nn.utils.clip_grad_norm_(model.parameters(), 1.0)
        optimizer.step()
    scheduler.step()
    print(f"Epoch {epoch}: loss={loss.item():.4f}")
\end{verbatim}
\end{LTR}

\section{סיכום}

פרק זה הציג את שיטות האימון והאופטימיזציה המרכזיות בלמידה עמוקה: מירידת
גרדיאנט בסיסית (\textenglish{SGD}) דרך אלגוריתמים מתקדמים כמו \textenglish{Adam}
ו-\textenglish{AdamW}, ועד לטכניקות רגולריזציה ותזמון קצב למידה. בפרק הבא נדון
ביישומים מעשיים ובהערכת מודלים.
